\subsection{Meta-RL} \label{sec:methods_problem} 

We use a black-box meta-RL problem setting \citep{duan2017rl, wang2016learning}. We define the task space $\mathcal{M}$ to be a set of partially-observable Markov decision processes (POMDPs). For a given task $m \in \mathcal{M}$ we define a \emph{trial} to be any sequence of transitions from an initial state $s_0$ to a terminal state $s_T$.\footnote{Note that we use a reversed naming convention to \citet{duan2017rl}. In our convention, the term ``trial'' maps well onto the related concept in the human behavioural literature \citep{barbosa2022practical}.} In XLand, tasks terminate if and only if a certain time period $T \in [10\textrm{s}, 40\textrm{s}]$ has elapsed, specified per-task. The environment ticks at $30$ frames-per-second and the agent observes every $4$\textsuperscript{th} frame, so task lengths in units of timesteps lie in the range $[75, 300]$.

An \emph{episode} consists of a sequence of $k$ trials for a given task $m$. At trial boundaries, the task is reset to an initial state. In our domain, initial states are deterministic except for the rotation of the agent, which is sampled uniformly at random. The trial and episode structure is depicted in Figure \ref{fig:xland-overview}.

In black-box meta-RL training, an agent uses experience of interacting with a wide distribution of tasks to update the parameters of its neural network, which parameterises the agent's policy distribution over actions given a state observation. If an agent possesses dynamic internal state (memory), then meta-RL training endows that memory with an implicit online learning algorithm, by leveraging the structure of repeated trials \citep{mikulik2020metatrained}. 

At test time, this online learning algorithm enables the agent to adapt its policy without any further updates to the neural network weights. Therefore, the memory of the agent is not reset at trial boundaries, but is reset at episode boundaries. To generate an episode, we sample a pair ($m$, $k$) where $k \in \{1,2, \dots 6\}$. As we will discuss later, at test time AdA is evaluated on unseen, held-out tasks across a variety of $k$ values, including on held-out $k$ not seen during training. For full details on AdA's meta-RL method, see Appendix \ref{appendix:metarl}. 